\begin{lemma}[Defective-chain structure, part (b)]
Let $(T,e_0,\varphi)$ be a globally maximum triple for $G$.  Let
$\alpha\in\mathcal D(V(T),\varphi)$ be defective, let
$\gamma\in\mathcal F(T,e_0,\varphi)$ be free and missing at a vertex
$u\in V(T)$, and let $z\notin V(T)$.  If the supplied incidence walk $P$
from $u$ to $z$ is a full defective-chain path for
$(T,\varphi,\alpha,\gamma,u,z)$, then
\[
 \partial_{\varphi,\alpha}(V(T))
 =\{f\in\partial_G(V(T)): \text{the edge-node of }f\text{ lies on }P\}.
\]
The equality is between sets of edge identities, including distinct
parallel edges.
\end{lemma}
\begin{proof}
Put $X=V(T)$, with $X$ understood as the vertex set supplied by
\noderef{TreeSequenceVertexSet}.  We first prove the right-to-left
inclusion.  Let $f\in\partial_G(X)$ and suppose that the edge-node of $f$
lies in the support of $P$.  The ends of $P$ are vertex-nodes.  Hence this
edge-node is incident along the walk with a vertex-node.  By
\noderef{TwoColorIncidenceGraph}, its color under $\varphi$ is therefore
either $\alpha$ or $\gamma$.  Since $\gamma$ is free,
\noderef{TreeFreeColors} and \noderef{TreeMissingColors} show that
$\gamma$ is missing at some vertex of $X$.  The missing-color clause of
\noderef{ScheideGlobalMaximumBoundaryFacts} says that no boundary edge of
$X$ has color $\gamma$.  Since $f$ is a boundary edge in the sense of
\noderef{BoundaryEdges}, its color must be $\alpha$.  The definition
\noderef{BoundaryColorEdges} now gives
$f\in\partial_{\varphi,\alpha}(X)$.

For the reverse inclusion, fix
$f\in\partial_{\varphi,\alpha}(X)$ and suppose, for a contradiction, that
the edge-node of $f$ does not lie on $P$.  By
\noderef{BoundaryColorEdges}, $f\in\partial_G(X)$ and
$\varphi(f)=\alpha$.  Let
\[
 Q=\noderef{TwoColorComponent}(G,\varphi;\alpha,\gamma,
       \operatorname{inr}f)
\]
be the complete two-color incidence component through that particular
edge identity.  The full-path clause of
\noderef{ScheideDefectiveChainPath} identifies the support of $P$ with
the complete component through the vertex-node $u$.  If $u$ belonged to
$Q$, then $u$ and the edge-node of $f$ would be connected, so that edge-node
would belong to the component through $u$ and hence to $P$, contrary to
our assumption.  Thus $u\notin Q$; in particular $Q$ is distinct from the
component displayed by $P$.

We record the two-color incidences at vertices of $X$.  The defective
boundary clause of \noderef{ScheideGlobalMaximumBoundaryFacts} supplies
at least one $\alpha$-colored boundary edge.  On the other hand,
\noderef{ScheideGlobalMaximumOrder} says that $X$ is closed.  Unfolding
closedness with \noderef{ColoredVertexSetClosed}, that boundary edge shows
that $\alpha$ is not missing at any $x\in X$.  Since
\noderef{MissingColors} is the complement of
\noderef{PresentColors}, $\alpha$ is present at every vertex of $X$.
The same maximum-order result says that $X$ is elementary.  By
\noderef{ColoredVertexSetElementary}, the fact that $\gamma$ is missing
at $u$ implies that it is not missing at any $x\in X\setminus\{u\}$;
hence \noderef{MissingColors} and \noderef{PresentColors} show that
$\gamma$ is present at each such $x$.

Every present color at a vertex has exactly one incident edge identity.
Existence is the defining witness in \noderef{PresentColors}.  For
uniqueness, two distinct edges of the same color incident with one vertex
would be adjacent by \noderef{Incident} and \noderef{AdjacentEdges}.  But
the globally maximum hypothesis in \noderef{GlobalMaximumTriple} includes
a \noderef{KTriple}, whose \noderef{IsProperEdgeColoring} clause forbids
two distinct adjacent colored edges from receiving the same color.  In
particular, at $u$ there is one $\alpha$-edge and no $\gamma$-edge, while
at every $x\in X\setminus\{u\}$ there is exactly one edge of each color.
This also proves $\alpha\ne\gamma$: $\alpha$ is present at $u$, whereas
$\gamma\in\noderef{MissingColors}(G,\varphi,u)$.

We next determine exactly what the switch on $Q$ does at the vertices of
$X$.  Every edge-node in $Q$ has original color $\alpha$ or $\gamma$.
Indeed, the node of $f$ has color $\alpha$ and is adjacent to its
endpoints; every other edge-node in its connected component occurs on a
nontrivial incidence walk, and \noderef{TwoColorIncidenceGraph} permits
an incidence adjacency at an edge-node only for those two colors.  If a
vertex-node $x\in X$ lies in $Q$, then $x\ne u$.  Its unique incident
$\alpha$- and $\gamma$-edges are adjacent to it in the incidence graph,
so both edge-nodes also lie in $Q$.  There are no other incident edge-nodes
of $Q$ at $x$, by the preceding uniqueness and two-color observations.
Thus switching $Q$ merely exchanges these two incident edges at $x$ and
leaves the set of colors present there unchanged.  If instead
$x\notin Q$, no edge-node of $Q$ is incident with $x$: such an incident
two-colored edge-node would be adjacent to $x$ in
\noderef{TwoColorIncidenceGraph} and would put $x$ in the same component.
Consequently the switch changes no incident edge at such a vertex.

Define the switched partial coloring explicitly by
\[
 \varphi'=\noderef{TwoColorComponentSwap}
   (G,\varphi;\alpha,\gamma,\operatorname{inr}f).
\]
The preceding paragraph and the definition of
\noderef{TwoColorComponentSwap} show, for every $x\in X$, that
\[
 \noderef{PresentColors}(G,\varphi',x)
   =\noderef{PresentColors}(G,\varphi,x).
\]
Taking complements via \noderef{MissingColors} gives equality of the
missing-color sets at every tree vertex.  In particular $\gamma$ remains
missing at $u$.  Since $\alpha\ne\gamma$, the properness statement
\noderef{TwoColorComponentSwapProper} applies and proves that $\varphi'$
is proper.

We also check the distinguished edge and all displayed tree edges.  The
$k$-triple clause of \noderef{GlobalMaximumTriple}, expanded through
\noderef{KTriple} and \noderef{IsOneEdgeUncolored}, says that
$\varphi(e_0)$ is uncolored and every other edge is colored.  The swap in
\noderef{TwoColorComponentSwap} fixes uncolored values, sends either of
the two switched colors to the other color, and fixes all other values.
Therefore $e_0$ is still uncolored under $\varphi'$, and every
$e\ne e_0$ is still colored; $e_0$ remains the unique uncolored edge.

Moreover, \noderef{ScheideGlobalMaximumUnused} says that both the
defective color $\alpha$ and the free color $\gamma$ are absent from the
used-color set \noderef{TreeUsedColors}.  Hence no displayed colored tree
edge has either color.  Since every edge-node in $Q$ has one of those two
colors, no displayed colored tree edge lies in $Q$; the first displayed
edge is $e_0$ and is uncolored, so it does not lie in $Q$ either.  Thus
every displayed tree-edge value is unchanged by the switch.

It follows clause by clause that the old sequence $T$ is still a
Tashkinov tree for $(e_0,\varphi')$.  Its vertices, edge identities,
distinctness conditions, attachments, and first-edge identity are the
same, so it is still the same \noderef{TreeSequence}.  For every displayed
edge after the first, the old \noderef{TashkinovTree} witness consists of
its color and an earlier vertex missing that color.  We have proved both
that the displayed edge keeps its color and that the full missing-color
set at every tree vertex is unchanged.  Hence every old witness remains
valid under $\varphi'$.

It remains to append $f$ and verify every new condition.  Since
$f\in\partial_G(X)$, the definition \noderef{BoundaryEdges} supplies
endpoints $x\in X$ and $y\notin X$ with
$\operatorname{ends}(f)=\{x,y\}$.  Write the ordered data of
\noderef{TreeSequence} as
\[
 T=(y_0,e_1,y_1,\ldots,e_n,y_n),
 \qquad e_1=e_0,
\]
and form the exact appended sequence
\[
 T'=(y_0,e_1,y_1,\ldots,e_n,y_n,f,y).
\]
The vertex $y$ is new because $y\notin X$, while $x\in X$ is one of the
earlier displayed vertices and so gives the required attachment of $f$.
The edge identity $f$ is new: its edge-node lies in $Q$, whereas we proved
that no displayed tree-edge node lies in $Q$.  All old vertices and edges
remain pairwise distinct, the new ones are distinct from them, every old
attachment is unchanged, and the first edge is still $e_1=e_0$.
Therefore $T'$ satisfies every clause of \noderef{TreeSequence}, and its
vertex set is $X\cup\{y\}$ with cardinality $|X|+1$.

By the definition of the switch, the chosen identity $f$, which originally
has color $\alpha$ and lies in $Q$, now has color $\gamma$.  The vertex
$u\in X$ occurs earlier than the appended vertex, and we proved that
$\gamma$ is still missing there.  Thus $u$ is the missing-color witness
required by \noderef{TashkinovTree} for the new last edge $f$; all earlier
edges retain the witnesses checked above.  Hence $T'$ is a Tashkinov tree
for $(e_0,\varphi')$.

Finally, the graph-criticality, chromatic-index equality, and maximum-degree
inequality in \noderef{KTriple} do not depend on the coloring.  Together
with the properness and unique-uncolored conclusions already proved, they
show that $(G,e_0,\varphi')$ is again a $k$-triple.  Its chromatic-index
equality gives $k=\chi'(G)-1$, so the definition
\noderef{TashkinovOrder} applies to $T'$ and yields
$|V(T')|\leq t(G)$.  But \noderef{GlobalMaximumTriple} gives
$|X|=t(G)$, whereas the append construction gives
$|V(T')|=|X|+1$.  This contradiction proves that the edge-node of every
$f\in\partial_{\varphi,\alpha}(X)$ lies on $P$.  Together with the first
inclusion, this proves the claimed equality of edge identities.
\end{proof}
